\section{Physical derivative and full vector-force costs of the five-row mean operation}
\label{mcmean:section}

This section completes the finite mean operation on the actual current
source state. All coordinates are the original right-handed cylindrical
coordinates. The physical viscosity is the fixed number
\(\nu_{\rm NS}>0\). Write the full current velocity as
\(U=(b+\beta,V+v,G+\gamma)+w\), where \(w\) includes every current
wave and curl correction. During the operation below this wave is held
fixed. The three targets are the full current defects
\begin{equation}\label{mcmean:full_targets}
 \mathcal P=P_{\rm old}+\Delta P,\qquad
 \mathcal J_\theta=(J_\theta)_{\rm old}+\Delta J_\theta,\qquad
 \mathcal J_z=(J_z)_{\rm old}+\Delta J_z.
\end{equation}
Here the increments are the pressure-corrected actual wave increments
of \eqref{mcq:defects}. If other finite source operations have already
been performed, ``old'' means their actual entering state and the
increments are recomputed for the actual addition. Thus none of the
following formulas substitutes an increment for a full target.

Retain, without a change of scale or profile constant,
\begin{equation}\label{mcmean:original_scale}
\begin{gathered}
 q-z^2q^{2h}=1-t,\quad z=q^D\eta,\quad x=r/\sqrt q,\quad
 A=\tfrac12+h,\quad D=\tfrac12-h,\\
 0<h<1/100,\quad q>0,\quad -1<\eta<1,\quad
 B(\eta)=1-2h\eta^2,\\
 V=q^{-A}a(\eta)x^{-1-2\lambda},\quad G=0
 \quad\hbox{on the reserved mean patch},\\
 a(\eta)=\frac{2^{1/2+\lambda}c_{\rm patch}}{1+\eta^2},
 \qquad \lambda>0,\qquad |a(\eta)|\ge a_0>0.
\end{gathered}
\end{equation}
The symbol \(a(\eta)\) in this display is a scalar base coefficient;
below the bold symbol \(\boldsymbol a\) denotes the vector increment.
The original constant \(c_{\rm patch}\) can have either nonzero sign.

\subsection{Exact inverse profiles and the five physical rows}

Use precisely the bumps and spacing of the five-row operation:
\[
 \varphi_j(x)=e^{-jd_b}\varphi_0(e^{-jd_b}x),\qquad
 \mu_p=\int_0^\infty x^p\varphi_0(x)\,dx>0,
 \qquad t_p=e^{d_bp},\quad d_b>0.
\]
The nonnegative nonzero smooth bump \(\varphi_0\) and its three
dilates have support in the reserved open interval \(I_m\Subset(0,\infty)\).
The substitution \(x=e^{jd_b}\xi\) proves exactly
\(\int x^p\varphi_j(x)\,dx=\mu_p t_p^j\).
Number matrix rows and columns starting at zero, and define
\begin{equation}\label{mcmean:vandermonde}
\begin{gathered}
 V_\theta=[t_p^j]_{p=2,-2-2\lambda,-2\lambda; j=0,1,2},
 \qquad B_\theta=V_\theta^{-1},\\
 V_z=[t_p^j]_{p=1,1-2\lambda; j=0,1},
 \qquad B_z=V_z^{-1},\\
 A_\theta=\operatorname{diag}
  (\mu_2,2a\mu_{-2-2\lambda},-a\mu_{-2\lambda})V_\theta,
 \qquad
 A_z=\operatorname{diag}(\mu_1,a\mu_{1-2\lambda})V_z.
\end{gathered}
\end{equation}
For three row values \(t_0,t_1,t_2\), subtraction of the first row
from each later row, followed by factoring, gives
\(\det V_\theta=(t_1-t_0)(t_2-t_0)(t_2-t_1)\).
For two row values the determinant is \(t_1-t_0\).
The original exponents in each displayed matrix are distinct because
\(\lambda>0\), and the exponential with \(d_b>0\) is injective.
The row factors are nonzero. This proves both inverses exist at the
actual fixed source parameters. Explicitly each entry of either
\(B\) is its corresponding transposed cofactor divided by the stated
determinant; this is a finite formula in the original exponentials.

Define the following actual profiles, including their \(\eta\) dependence:
\begin{equation}\label{mcmean:profiles}
\begin{split}
 V_P(x,\eta)&=-\frac{1}{2a(\eta)\mu_{-2-2\lambda}}
                   \sum_{j=0}^2B_\theta[j,1]\varphi_j(x),\\
 V_{J_z}(x,\eta)&=\frac{1}{a(\eta)\mu_{-2\lambda}}
                   \sum_{j=0}^2B_\theta[j,2]\varphi_j(x),\\
 G_{J_\theta}(x,\eta)&=-\frac{1}{a(\eta)\mu_{1-2\lambda}}
                   \sum_{j=0}^1B_z[j,1]\varphi_j(x).
\end{split}
\end{equation}
The signs, the factor two, and all three negative powers in the moment
denominators are part of the formulas. Multiplying the columns by
the exact rows in \eqref{mcmean:vandermonde} proves the eight scalar
profile identities
\begin{equation}\label{mcmean:profile_moments}
\begin{array}{c|cc}
 &V_P&V_{J_z}\\ \hline
 \int x^2(\cdot)\,dx&0&0\\
 2a\int x^{-2-2\lambda}(\cdot)\,dx&-1&0\\
 -a\int x^{-2\lambda}(\cdot)\,dx&0&-1
\end{array}
\qquad
\begin{array}{c|c}
 &G_{J_\theta}\\ \hline
 \int x(\cdot)\,dx&0\\
 a\int x^{1-2\lambda}(\cdot)\,dx&-1
\end{array}.
\end{equation}
There are eight entries in these two tables; each follows from
\(V_\theta B_\theta=I_3\) or \(V_zB_z=I_2\), with the row factor
left in place. The complete matrix identities include all nine and
all four scalar inverse identities, respectively.

The source coefficients are
\[
 \boldsymbol u=A_\theta^{-1}
       (0,-q^{2A}\mathcal P,-q^{2A-1}\mathcal J_z)^T,
 \qquad
 \boldsymbol s=A_z^{-1}
       (0,-q^{2A-3/2}\mathcal J_\theta)^T.
\]
Since \((DV)^{-1}=V^{-1}D^{-1}\), substitution of these coefficients
in the original physical increments gives
\begin{equation}\label{mcmean:tangential_increment}
 \delta v=q^A\mathcal P V_P+q^{A-1}\mathcal J_zV_{J_z},
 \qquad
 \delta\gamma=q^{A-3/2}\mathcal J_\theta G_{J_\theta}.
\end{equation}
For example the fifth-row coefficient in \(\boldsymbol u\) is
\((-q^{2A-1}\mathcal J_z)/(-a\mu_{-2\lambda})\), which proves
the positive sign in \(V_{J_z}\). Multiplication by the original
outer factor \(q^{-A}\) proves its exponent \(A-1\).
The other two terms follow by the same displayed diagonal products.

Changing variables \(r=\sqrt q\,x\) in each integral, and using
\eqref{mcmean:profile_moments}, now proves the full five equations
\begin{equation}\label{mcmean:five_rows}
\begin{gathered}
 \int r^2\delta v\,dr=0,\qquad
 \int r\delta\gamma\,dr=0,\qquad
 \int\frac{2V}{r}\delta v\,dr=-\mathcal P,\\
 \int r^2(G\delta v+V\delta\gamma)\,dr=-\mathcal J_\theta,\\
 \int(2rG\delta\gamma-rV\delta v)\,dr=-\mathcal J_z.
\end{gathered}
\end{equation}
Indeed the pressure term with \(V_P\) has scale
\(q^{-A}q^A=1\); the angular flux with \(G_{J_\theta}\) has scale
\(q^{3/2}q^{-A}q^{A-3/2}=1\); and the fifth row with
\(V_{J_z}\) has scale \(q q^{-A}q^{A-1}=1\).
Every other term is one of the zero entries in the tables. The two
homogeneous velocity moments use the first rows of those tables.
This verifies the physical equations, including the single factor
\(r\) multiplying \(V\delta v\) in the fifth row.

\subsection{Compact potential, induced radial velocity, and support}

Define the combined potential profile
\begin{equation}\label{mcmean:H}
 H_\theta(x,\eta)=\frac1x\int_0^x\xi G_{J_\theta}(\xi,\eta)\,d\xi.
\end{equation}
Choose \(0<x_-<x_+\) such that the convex hull of all the bump
supports lies in \([x_-,x_+]\Subset I_m\). For \(x<x_-\) the
integral in \eqref{mcmean:H} is zero. For \(x>x_+\) it is the
zero axial moment in \eqref{mcmean:profile_moments}.
It follows that \(H_\theta\), together with every derivative, is
supported in this fixed compact interval. Its smoothness at the two
collars follows from the fact that the integrand is identically zero
there and that its total integral is zero. In particular it extends
smoothly by zero to \(x\le0\).

The primitive of an individual \(\varphi_j\) is
\(\psi_j=x^{-1}\int_0^x\xi\varphi_j(\xi)\,d\xi\), which equals
\(\mu_1e^{jd_b}/x\) past its support. Thus an individual primitive
can have a tail. The exact identity
\(\sum_{j=0}^1 B_z[j,1]e^{jd_b}=0\) cancels these tails in
\(H_\theta=-(a\mu_{1-2\lambda})^{-1}\sum_j B_z[j,1]\psi_j\).
The identity holds for every \(\eta\), so its \(\eta\) derivatives
also vanish. This proves compactness for the combined profile and
all derivatives without assuming compactness for the summands.

The physical potential and radial increment are exactly
\begin{equation}\label{mcmean:potential_radial}
\begin{split}
 \Psi(r,z,t)&=\frac1r\int_0^r s\delta\gamma(s,z,t)\,ds
             =q^{A-1}\mathcal J_\theta H_\theta(x,\eta),\\
 \delta\beta&=-\partial_z\Psi\\
 &=-q^{A-1}(\partial_z\mathcal J_\theta)H_\theta
   -q^{A-1-D}\mathcal J_\theta\mathcal Z_{A-1}H_\theta.
\end{split}
\end{equation}
The substitution in the integral supplies one factor \(q\) and its
outer \(1/r\) supplies \(q^{-1/2}\), giving the exponent \(A-1\).
The operator in the second line is explicitly
\[
 \mathcal Z_\sigma f
 =\frac{2\sigma\eta f-\eta x f_x+(1-\eta^2)f_\eta}{B(\eta)}.
\]
The chain rule \eqref{mcq:dilation_dictionary} proves the last line,
including the \(\eta\) derivative of the original factor \(1/a\).
Direct differentiation of \(xH_\theta\) gives
\((\partial_x+x^{-1})H_\theta=G_{J_\theta}\), so
\((\partial_r+r^{-1})\Psi=\delta\gamma\). Commuting the smooth
physical \(r,z\) derivatives proves
\begin{equation}\label{mcmean:divergence}
 (\partial_r+r^{-1})\delta\beta+\partial_z\delta\gamma=0.
\end{equation}
The three velocity increments and \(\Psi\) are auxiliary independent
and angularly independent; their support lies in
\(\sqrt q[x_-,x_+]\), including any intervals between the bumps.

\subsection{Every mixed physical derivative, with original scale powers}

For clarity the original differential dictionary is
\[
 q_z=\frac{2\eta q^A}{B},\quad q_t=-\frac1B,\quad
 \eta_z=\frac{q^{-D}(1-\eta^2)}B,\quad
 \eta_t=\frac{D\eta q^{-1}}B,
 \qquad
 \mathcal T_\sigma f=
 \frac{-\sigma f+\tfrac12x f_x+D\eta f_\eta}{B}.
\]
These equations follow by implicit differentiation of
\eqref{mcmean:original_scale}; in particular
\(1-2hz^2q^{2h-1}=B\) and \(D+2h=A\).
Differentiating \(x=rq^{-1/2}\) supplies
\(x_z=-\eta xq^{-D}/B\) and \(x_t=xq^{-1}/(2B)\).
Consequently
\[
 \partial_r(q^\sigma F)=q^{\sigma-1/2}F_x,\quad
 \partial_z(q^\sigma F)=q^{\sigma-D}\mathcal Z_\sigma F,\quad
 \partial_t(q^\sigma F)=q^{\sigma-1}\mathcal T_\sigma F.
\]
Here \(F=F(x,\eta)\). These are physical derivatives of the moving
scale; the time loss is the original one in this scale.

For \(k,m,n\ge0\) define \(F_{\sigma;k,m,n}\) by the following
finite recursion. Start with \(f=\partial_x^kF\) and
\(s=\sigma-k/2\). Repeat
\(f\leftarrow\mathcal Z_s f,\ s\leftarrow s-D\) exactly \(m\)
times; then repeat \(f\leftarrow\mathcal T_s f,\ s\leftarrow s-1\)
exactly \(n\) times. The final \(f\) is \(F_{\sigma;k,m,n}\).
For any actual scalar target \(X=X(z,t)\), the exact product formula is
\begin{equation}\label{mcmean:all_derivatives}
\begin{split}
 \partial_r^k\partial_z^m\partial_t^n(q^\sigma XF)
  =\sum_{i=0}^m\sum_{j=0}^n&\binom mi\binom nj
     (\partial_z^i\partial_t^jX)\\
  &\cdot q^{\sigma-k/2-(m-i)D-(n-j)}
       F_{\sigma;k,m-i,n-j}(x,\eta).
\end{split}
\end{equation}
To prove it, first apply \(k\) radial derivatives, which do not act
on \(X,q,\eta\). The repeated product rule for the \(m\) axial and
\(n\) time derivatives assigns \(i,j\) derivatives to \(X\), with
the two binomial multiplicities displayed. Apply the three chain
rules above successively to the remaining factor. Each axial
derivative lowers its current exponent by \(D\), and each time
derivative lowers it by one, while its current exponent is the index
in the next operator. This is precisely the stated recursion.
Smooth mixed derivatives commute, so the chosen radial-then-axial-
then-time order computes the derivative written in the formula.

This proves every mixed derivative of \(\delta v,\delta\gamma,\Psi\)
by using, respectively, the two terms and the one term in
\eqref{mcmean:tangential_increment} and \eqref{mcmean:potential_radial}.
It proves every derivative of \(\delta\beta\) either by applying
\eqref{mcmean:all_derivatives} to its two displayed terms or by
\begin{equation}\label{mcmean:radial_all_derivatives}
 \partial_r^k\partial_z^m\partial_t^n\delta\beta
 =-\partial_r^k\partial_z^{m+1}\partial_t^n
          (q^{A-1}\mathcal J_\theta H_\theta).
\end{equation}
This equality retains, in particular, the derivative of the full
target and the additional original factor \(q^{-D}\).

The recursion is an explicit finite rational expression in
\(x,\eta,h\), the retained profile derivatives, and powers of
\(B^{-1}\). The other profile constants are exactly the inverse
Vandermonde entries and moments in \eqref{mcmean:profiles}.
In fact the reciprocal base coefficient has the exact derivatives
\begin{equation}\label{mcmean:reciprocal_a}
 \partial_\eta^s(a^{-1})=
 \frac{1}{2^{1/2+\lambda}c_{\rm patch}}
 \begin{cases}1+\eta^2&s=0,\\2\eta&s=1,\\2&s=2,\\0&s\ge3.
 \end{cases}
\end{equation}
Also \(\varphi_j^{(k)}(x)=e^{-(k+1)jd_b}
\varphi_0^{(k)}(e^{-jd_b}x)\).
These two formulas give every mixed derivative of the three profiles
without an unspecified inverse-matrix derivative constant.

For a similarly explicit derivative of \(H_\theta\), write
\(F_0(x,\eta)=\int_0^x\xi G_{J_\theta}(\xi,\eta)\,d\xi\).
For \(\ell\ge1\) the fundamental theorem of calculus and the
product rule give
\[
 \partial_x^\ell F_0=x\partial_x^{\ell-1}G_{J_\theta}
              +(\ell-1)\partial_x^{\ell-2}G_{J_\theta},
\]
where the second term is defined to be zero for \(\ell=1\).
Therefore, for \(x>0\),
\begin{equation}\label{mcmean:H_derivatives}
 \partial_x^k\partial_\eta^sH_\theta
 =\sum_{\ell=0}^k\binom k\ell(-1)^{k-\ell}(k-\ell)!
    x^{-1-k+\ell}\partial_x^\ell\partial_\eta^sF_0.
\end{equation}
The \(\ell=0\) factor is the explicitly specified integral; every
other factor is the displayed profile derivative. For example its
absolute integral is at most
\(\int_{x_-}^{x_+}\xi|\partial_\eta^sG_{J_\theta}(\xi,\eta)|\,d\xi\).
This supplies finite computable profile costs on the exact support.
The bound \(B\ge1-2h>0\), together with \(x\ge x_->0\), proves
that each fixed recursively specified coefficient is bounded on
\([x_-,x_+]\times[-1,1]\). The original moments, spacing, \(h\),
\(c_{\rm patch}\), and \(\lambda\) remain in those coefficients.

\subsection{Explicit costs in full target derivative arrays}

For \(I=(m,n)\), put \(I_z=(m+1,n)\). Use the original actual-shape
arrays \(\mathcal B_{e,I}\) from \eqref{mcq:B}, with their two
retained coupled amplitude components and old-wave cross terms.
Define the full target majorants
\begin{equation}\label{mcmean:target_arrays}
\begin{split}
 \mathsf P_I&=|\partial_z^m\partial_t^nP_{\rm old}|
          +\mathcal B_{-1,I}+\tfrac12\mathcal B_{0,I_z},\\
 \mathsf J_{\theta,I}
      &=|\partial_z^m\partial_t^n(J_\theta)_{\rm old}|
          +\tfrac12\mathcal B_{2,I},\\
 \mathsf J_{z,I}
      &=|\partial_z^m\partial_t^n(J_z)_{\rm old}|
          +\mathcal B_{1,I}+\tfrac14\mathcal B_{2,I_z}.
\end{split}
\end{equation}
The triangle inequality in \eqref{mcmean:full_targets} and the
proved increment bounds \eqref{mcq:defect_bounds} give
\(|\partial^I\mathcal P|\le\mathsf P_I\),
\(|\partial^I\mathcal J_\theta|\le\mathsf J_{\theta,I}\), and
\(|\partial^I\mathcal J_z|\le\mathsf J_{z,I}\).
Thus the old absolute derivatives appear explicitly in every cost.

For a nonnegative target array \(\mathsf X\), define the following
pointwise, fully specified finite sum:
\begin{equation}\label{mcmean:cost_operator}
\begin{split}
 \mathfrak C_{\sigma,F;k,m,n}[\mathsf X](r,z,t)
  =\sum_{i=0}^m\sum_{j=0}^n&\binom mi\binom nj\mathsf X_{(i,j)}(z,t)\\
 &\cdot q^{\sigma-k/2-(m-i)D-(n-j)}
       |F_{\sigma;k,m-i,n-j}(x,\eta)|.
\end{split}
\end{equation}
All powers are powers of the actual positive \(q\), even when their
exponents are negative. Define
\begin{equation}\label{mcmean:velocity_arrays}
\begin{split}
 \mathsf A_{\theta;k,m,n}
    &=\mathfrak C_{A,V_P;k,m,n}[\mathsf P]
        +\mathfrak C_{A-1,V_{J_z};k,m,n}[\mathsf J_z],\\
 \mathsf A_{z;k,m,n}
    &=\mathfrak C_{A-3/2,G_{J_\theta};k,m,n}[\mathsf J_\theta],\\
 \mathsf S_{k,m,n}
    &=\mathfrak C_{A-1,H_\theta;k,m,n}[\mathsf J_\theta],\\
 \mathsf A_{r;k,m,n}
    &=\mathfrak C_{A-1,H_\theta;k,m,n}[\mathsf J_\theta^{+z}]
      +\mathfrak C_{A-1-D,\mathcal Z_{A-1}H_\theta;k,m,n}
                                                 [\mathsf J_\theta],
 \qquad (\mathsf J_\theta^{+z})_{(i,j)}=\mathsf J_{\theta,(i+1,j)}.
\end{split}
\end{equation}
Equations \eqref{mcmean:all_derivatives}--\eqref{mcmean:potential_radial}
and the triangle inequality prove, for every \(k,m,n\),
\begin{equation}\label{mcmean:velocity_bounds}
\begin{gathered}
 |\partial_r^k\partial_z^m\partial_t^n\delta v|
       \le\mathsf A_{\theta;k,m,n},\qquad
 |\partial_r^k\partial_z^m\partial_t^n\delta\gamma|
       \le\mathsf A_{z;k,m,n},\\
 |\partial_r^k\partial_z^m\partial_t^n\Psi|
       \le\mathsf S_{k,m,n},\qquad
 |\partial_r^k\partial_z^m\partial_t^n\delta\beta|
       \le\mathsf A_{r;k,m,n}.
\end{gathered}
\end{equation}
The alternative upper bound \(\mathsf S_{k,m+1,n}\) for the last
quantity follows directly from \eqref{mcmean:radial_all_derivatives}.
For example the zero-order costs read
\[
\begin{split}
 |\delta v|&\le q^A\mathsf P_{0,0}|V_P|
                    +q^{A-1}\mathsf J_{z,0,0}|V_{J_z}|,\\
 |\delta\gamma|&\le q^{A-3/2}\mathsf J_{\theta,0,0}|G_{J_\theta}|,
 \qquad |\Psi|\le q^{A-1}\mathsf J_{\theta,0,0}|H_\theta|,\\
 |\delta\beta|&\le q^{A-1}\mathsf J_{\theta,1,0}|H_\theta|
       +q^{A-1-D}\mathsf J_{\theta,0,0}|\mathcal Z_{A-1}H_\theta|.
\end{split}
\]
These display separately the target derivative and moving-scale radial
cost. Replacing each profile value in \eqref{mcmean:cost_operator}
by its maximum over its fixed \(x,\eta\) support supplies a pointwise
bound depending only on \(z,t\) and the stated fixed profile data.

The same formula gives weighted radial costs without a loss hidden in
a constant. For any real \(e\), replace the final profile absolute
value in \eqref{mcmean:cost_operator} by
\[
 q^{(e+1)/4}
 \left(\int_{x_-}^{x_+}x^e
          |F_{\sigma;k,m-i,n-j}(x,\eta)|^2\,dx\right)^{1/2}.
\]
Call the resulting sum \(\mathfrak C^{[e]}\). The exact measure
identity \(r^e\,dr=q^{(e+1)/2}x^e\,dx\), followed by the triangle
inequality in this weighted \(L^2\) space, proves all four bounds
\eqref{mcmean:velocity_bounds} with the left sides replaced by their
\([e]\) norms and the right sides by the corresponding
\(\mathfrak C^{[e]}\) sums. The support is bounded away from zero,
so these integrals are finite for every fixed real \(e\).

\subsection{Actual pressure reconstruction and the surviving defects}

Set \(\boldsymbol a=(a_r,a_\theta,a_z)
=(\delta\beta,\delta v,\delta\gamma)\). Before pressure is
changed, let \(\delta g_r\) denote the exact change in the source's
angularly averaged negative radial residual. The angular average here
retains auxiliary dependence. The actual pressure-defect change is
\begin{equation}\label{mcmean:pressure_defect_step}
 \delta P_{\rm step}=\int_0^\infty\overline{\delta g_r}\,dr
                    =P_{\rm new}-\mathcal P.
\end{equation}
This is the defect change caused by the present mean operation;
it is different from \(\Delta P\) in \eqref{mcmean:full_targets}.
The source pressure bump is still
\(\rho=q^{-1/2}\widehat\rho(x)\),
\(\int\widehat\rho\,dx=1\). The actual pressure increment for this
operation is
\begin{equation}\label{mcmean:actual_pressure}
 \delta p=\mathcal I_c f,\qquad
 f=\delta g_r-\rho\,\delta P_{\rm step}.
\end{equation}
The original physical shifted primitive is
\[
\begin{split}
 \mathcal If(r,Y)&=\int_0^r
  f(s,Y+(s^{d_r}-r^{d_r})v_r,z,t)\,ds,\\
 \mathcal Jf(r,Y)&=\int_0^\infty
  f(s,Y+(s^{d_r}-r^{d_r})v_r,z,t)\,ds,\\
 \mathcal I_cf&=\mathcal If-\chi_m\mathcal Jf,\\
 d_r&=2((1+h)\rho_g-h\kappa_s),\quad
 \rho_g=\frac{\log(4-\sqrt2)}{\log(4+\sqrt2)},\quad
 \kappa_s=10^{-5},\quad v_r=(1,-(\sqrt2-1)).
\end{split}
\]
Here \(\chi_m\) is the actual inner-zero, outer-one source cutoff;
the two suppressed arguments of each integrand are exactly \(z,t\).
The pressure entering the physical force is the physical evaluation of
this source function, with its original phase evaluation and all its
derivatives. Thus it includes any nonzero auxiliary modes of the
source reconstruction.

To verify the reconstruction identity, the source physical radial
operator \(\mathcal D_r\) cancels the derivative of the term
\(-r^{d_r}v_r\) in the shift. Endpoint differentiation gives
\(\mathcal D_r\mathcal If=f\) and
\(\mathcal D_r\mathcal Jf=0\), and hence
\begin{equation}\label{mcmean:pressure_remainder}
 \mathcal D_r\delta p-\delta g_r
 =-\rho\,\delta P_{\rm step}
                  -(\partial_r\chi_m)\mathcal Jf.
\end{equation}
Haar translation invariance and \eqref{mcmean:pressure_defect_step}
give \(\overline{\mathcal Jf}=0\), but this does not set
\(\mathcal Jf\) itself to zero. In the inner collar \(\mathcal If=0\)
and \(\chi_m=0\); in the outer collar \(\mathcal If=\mathcal Jf\)
and \(\chi_m=1\). This proves compact support of the actual pressure
and all derivatives. In particular its gradient is retained in the
force calculation below.

For completeness the exact remaining source terms can be written
entirely in the original current mean fields. With
\(\Delta_0=\mathcal D_r^2+r^{-1}\mathcal D_r+\partial_z^2\),
direct expansion of the radial conservative source gives
\begin{equation}\label{mcmean:Rg}
\begin{split}
 R_g:={}&\delta g_r-\frac{2V}{r}\delta v\\
 ={}&-\partial_t\delta\beta
 -(\mathcal D_r+r^{-1})
       \bigl(2b\delta\beta+2\beta\delta\beta+(\delta\beta)^2\bigr)\\
 &-\partial_z\bigl(b\delta\gamma+G\delta\beta+
       \beta\delta\gamma+\gamma\delta\beta+
       \delta\beta\delta\gamma\bigr)\\
 &+r^{-1}\bigl(2v\delta v+(\delta v)^2\bigr)
       +\nu_{\rm NS}(\Delta_0-r^{-2})\delta\beta.
\end{split}
\end{equation}
Indeed the radial conservative quadratic expression before the change
is \(-(\mathcal D_r+r^{-1})(b+\beta)^2
-\partial_z((b+\beta)(G+\gamma))+(V+v)^2/r\), together with
the wave covariance, time term, and viscosity. Subtract it from its
value after the mean change. The unchanged wave covariance cancels;
each remaining quadratic product expands to the terms displayed.
The term \(2V\delta v/r\) has been written separately in the
definition of \(R_g\). The ordinary time derivative on the increment
is exact because it is auxiliary independent.

The five physical rows then give the exact new defects
\begin{equation}\label{mcmean:nonlinear_defects}
\begin{split}
 P_{\rm new}&=\int\overline{R_g}\,dr,\\
 (J_\theta)_{\rm new}
   &=\int r^2\overline{\gamma\delta v+v\delta\gamma+
                                    \delta\gamma\delta v}\,dr,\\
 (J_z)_{\rm new}
   &=\int r\overline{2\gamma\delta\gamma+(\delta\gamma)^2}\,dr
                        -\tfrac12\int r^2\overline{R_g}\,dr.
\end{split}
\end{equation}
For the first line add \(\mathcal P\) to
\(\int(2V\delta v/r+\overline{R_g})\,dr\) and use the third
row. For the second, expanding \((G+\gamma+\delta\gamma)
(V+v+\delta v)-(G+\gamma)(V+v)\) gives base-linear terms
\(G\delta v+V\delta\gamma\), cancelled by the fourth row, and
exactly the three remaining terms displayed. For the third, expand
\((G+\gamma+\delta\gamma)^2-(G+\gamma)^2\) and the radial-source
moment \(-\tfrac12\int r^2\delta g_r\,dr\); the base-linear
terms are the fifth row. This proves all three equalities.
Neither their right sides nor the cutoff term in
\eqref{mcmean:pressure_remainder} has been removed.

\subsection{Exact full cylindrical physical force increment}

For the actual physical fields, put
\(\mathcal R(U,p)=\partial_tU+(U\cdot\nabla)U
-\nu_{\rm NS}\Delta U+\nabla p\). Ordinary derivatives in this
formula are derivatives after the physical source evaluation; they
therefore include all phase, frame, cutoff and scale derivatives of
the actual full \(U\). Expanding the two quadratic products proves
the vector identity
\begin{equation}\label{mcmean:vector_force}
 \delta\mathcal R=\partial_t\boldsymbol a
 +(U\cdot\nabla)\boldsymbol a+(\boldsymbol a\cdot\nabla)U
 +(\boldsymbol a\cdot\nabla)\boldsymbol a
 -\nu_{\rm NS}\Delta\boldsymbol a+\nabla\delta p.
\end{equation}
The pressure here is precisely \eqref{mcmean:actual_pressure}.

To derive every cylindrical component, the basis derivatives are
\(\partial_\theta e_r=e_\theta\),
\(\partial_\theta e_\theta=-e_r\), and
\(\partial_\theta e_z=0\). Thus for two arbitrary vector fields
\(X,Y\) the component vector of \((X\cdot\nabla)Y\) is
\[
 X_r\partial_rY+\frac{X_\theta}{r}\partial_\theta Y
 +X_z\partial_zY+\frac{X_\theta}{r}(-Y_\theta,Y_r,0).
\]
Apply this equality to \((U,\boldsymbol a)\),
\((\boldsymbol a,U)\), and \((\boldsymbol a,\boldsymbol a)\).
The scalar components of \(\boldsymbol a\) have zero angular
derivative, whereas those of the actual \(U\) generally do not.
Writing \(L_0=\partial_r^2+r^{-1}\partial_r+\partial_z^2\), the
complete result is
\begin{equation}\label{mcmean:force_r}
\begin{split}
 \delta\mathcal R_r={}&\partial_ta_r+U_r\partial_ra_r+U_z\partial_za_r
       +a_r\partial_rU_r+a_z\partial_zU_r
       +\frac{a_\theta}{r}\partial_\theta U_r\\
 &+a_r\partial_ra_r+a_z\partial_za_r
       -\frac{2U_\theta a_\theta+a_\theta^2}{r}
       -\nu_{\rm NS}(L_0-r^{-2})a_r+\partial_r\delta p,
\end{split}
\end{equation}
\begin{equation}\label{mcmean:force_theta}
\begin{split}
 \delta\mathcal R_\theta={}&\partial_ta_\theta
       +U_r\partial_ra_\theta+U_z\partial_za_\theta
       +a_r\partial_rU_\theta+a_z\partial_zU_\theta
       +\frac{a_\theta}{r}\partial_\theta U_\theta\\
 &+a_r\partial_ra_\theta+a_z\partial_za_\theta
       +\frac{U_ra_\theta+a_rU_\theta+a_ra_\theta}{r}
       -\nu_{\rm NS}(L_0-r^{-2})a_\theta
       +\frac1r\partial_\theta\delta p,
\end{split}
\end{equation}
\begin{equation}\label{mcmean:force_z}
\begin{split}
 \delta\mathcal R_z={}&\partial_ta_z+U_r\partial_ra_z+U_z\partial_za_z
       +a_r\partial_rU_z+a_z\partial_zU_z
       +\frac{a_\theta}{r}\partial_\theta U_z\\
 &+a_r\partial_ra_z+a_z\partial_za_z
       -\nu_{\rm NS}L_0a_z+\partial_z\delta p.
\end{split}
\end{equation}
The self terms in the radial equation include
\(-a_\theta^2/r\); those in the angular equation include
\(a_ra_\theta/r\). The old-wave cross terms occur in the actual
\(U\) and in its angular derivatives in all three equations.
The mean pressure is angularly invariant, so its angular derivative
is zero for this source operation; its radial and axial derivatives
are the actual reconstructed ones and are generally nonzero.

For the viscosity formula, applying the scalar cylindrical Laplacian
to the moving basis yields, for a general vector field \(X\),
\[
 \Delta X=
 \left(\Delta_sX_r-r^{-2}X_r-2r^{-2}\partial_\theta X_\theta,
       \Delta_sX_\theta-r^{-2}X_\theta+2r^{-2}\partial_\theta X_r,
       \Delta_sX_z\right),
 \quad\Delta_s=L_0+r^{-2}\partial_\theta^2.
\]
This follows by differentiating the two basis identities twice and
applying the product rule. Substituting the angularly independent
components of \(\boldsymbol a\) gives exactly
\(((L_0-r^{-2})a_r,(L_0-r^{-2})a_\theta,L_0a_z)\), proving all
three displayed viscous terms.
The radial expression before the pressure term, when angularly
averaged in the extended source domain and negated, is precisely
\(\delta g_r\): this follows from the definition of the negative
radial source residual and linearity of the average. It is also an
explicit formula for the input to \eqref{mcmean:actual_pressure} in
terms of the same full velocity and mean increments.

\subsection{Finite all-order force costs with every inverse-radius term}

Here are explicit derivative bounds for all three components, including
angular derivatives of the actual current velocity. Let
\(\alpha=(k,\ell,m,n)\in\mathbb N^4\) index the scalar component
derivative \(\partial^\alpha=\partial_r^k\partial_\theta^\ell
\partial_z^m\partial_t^n\), and let
\(e_r,e_\theta,e_z,e_t\) be the four coordinate multi-indices.
Extend the arrays in \eqref{mcmean:velocity_arrays} by
\[
 \mathsf A_{i,\alpha}=
 \begin{cases}\mathsf A_{i;k,m,n}&\ell=0,\\0&\ell>0,
 \end{cases}
 \qquad
 \mathsf U_{i,\alpha}=|\partial^\alpha U_i|.
\]
The pressure array \(\mathsf\Pi\) will be the following explicit
majorant for the derivatives of the actual specified function
\eqref{mcmean:actual_pressure}, after physical evaluation.

The original physical radial phase satisfies
\(Y_{\rm phys}(s,z,t)-Y_{\rm phys}(r,z,t)
=(s^{d_r}-r^{d_r})v_r\). Substituting this exact equality in the
two shifted integrals proves the evaluated formula
\begin{equation}\label{mcmean:evaluated_pressure}
\begin{split}
 \delta p(r,z,t)
 &=\int_0^r f_{\rm phys}(s,z,t)\,ds
        -\chi_m(r,z,t)\int_0^\infty f_{\rm phys}(s,z,t)\,ds,\\
 f_{\rm phys}&=(\delta g_r)_{\rm phys}
                       -\rho\,\delta P_{\rm step}.
\end{split}
\end{equation}
Every original axial and temporal phase dependence is retained in
\(f_{\rm phys}\). The second integral is independent of \(r\)
after this evaluation. This equality does not assert that it is zero.

Fix any compact interior physical domain for the finite operation,
and choose one physical interval \([R_-,R_+]\Subset(0,\infty)\)
containing the supports of \(\delta g_r\) and \(\rho\) for that
domain and every auxiliary phase. Put \(L_R=R_+-R_-\), with these
actual endpoints kept fixed during differentiation. The smooth zero
extensions imply that all their derivatives have support in the
same interval. Define \(\mathsf G_{k,m,n}(z,t)\) as the supremum
over \(r\in[R_-,R_+]\) and the original auxiliary phases of the
angular average of the following nonnegative expression: the radial
cost \(\mathsf F_{r;(k,0,m,n)}\) in \eqref{mcmean:force_costs}
with its last pressure term deleted. This definition involves only
\(\mathsf A\), the actual current velocity derivatives,
\(\nu_{\rm NS}\), and the inverse-radius arrays; it has no
dependence on \(\mathsf\Pi\). When auxiliary phases are retained,
the velocity derivative arrays mean the corresponding full lifted
physical derivatives before evaluation. The chain rule for source
evaluation identifies these with derivatives of the evaluated fields
at each phase. The force identities and the derivative proof below
therefore imply
\[
 |\partial_r^k\partial_z^m\partial_t^n(\delta g_r)_{\rm phys}|
       \le\mathsf G_{k,m,n},\qquad
 |\partial_z^m\partial_t^n\delta P_{\rm step}|
       \le\mathsf D_{m,n}:=L_R\mathsf G_{0,m,n}.
\]
For the first inequality use the negative angularly averaged radial
force before pressure, as proved above, and the triangle inequality
under that average. For the second, differentiate
\eqref{mcmean:pressure_defect_step} under its compact radial
integral, commute axial and temporal differentiation with Haar
average, and integrate the same bound over the interval of width
\(L_R\). The Haar integral of each original constant torus
direction derivative is zero by periodicity, which justifies that
commutation for the full physical time derivative as well.

The derivative of the pressure bump is already an explicit original-
scale coefficient in \eqref{mcq:bump_derivatives}. Define its exact
support supremum
\[
 \mathsf H_{k,m,n}(z,t)
 =q^{-(1+k)/2-mD-n}
    \sup_{r\in[R_-,R_+]}
       |f_{0;k,m,n}(r/\sqrt q,\eta)|,
\]
where \(f_{0;k,m,n}\) is the recursion starting with the actual
\(\widehat\rho\), rather than another normalized bump.
It follows that \(|\partial_r^k\partial_z^m\partial_t^n\rho|
\le\mathsf H_{k,m,n}\). Set
\begin{equation}\label{mcmean:pressure_source_cost}
 \mathsf E_{k,m,n}
 =\mathsf G_{k,m,n}
 +\sum_{i=0}^m\sum_{j=0}^n\binom mi\binom nj
        \mathsf H_{k,m-i,n-j}\mathsf D_{i,j}.
\end{equation}
Since \(\delta P_{\rm step}\) has no radial dependence, all
\(k\) radial derivatives in its product with \(\rho\) act on
\(\rho\). The repeated axial and temporal product rule proves
\(|\partial_r^k\partial_z^m\partial_t^n f_{\rm phys}|
\le\mathsf E_{k,m,n}\) on the full interval.

Write \(\mathsf X_{k,i,j}(r,z,t)
=|\partial_r^k\partial_z^i\partial_t^j\chi_m(r,z,t)|\), retaining
the actual cutoff and its moving-scale derivatives. Define
\begin{equation}\label{mcmean:explicit_pressure_cost}
\begin{split}
 \mathsf\Pi_{(k,0,m,n)}={}&
 \begin{cases}L_R\mathsf E_{0,m,n}&k=0,\\
               \mathsf E_{k-1,m,n}&k\ge1
 \end{cases}\\
 &+L_R\sum_{i=0}^m\sum_{j=0}^n\binom mi\binom nj
           \mathsf X_{k,i,j}\mathsf E_{0,m-i,n-j},
 \qquad \mathsf\Pi_{(k,\ell,m,n)}=0\quad(\ell>0).
\end{split}
\end{equation}
This proves the explicit bound
\(|\partial^\alpha\delta p|\le\mathsf\Pi_\alpha\).
In fact differentiating \eqref{mcmean:evaluated_pressure} yields
the exact first-integral derivative
\[
 \partial_r^k\partial_z^m\partial_t^n
       \int_0^r f_{\rm phys}(s,z,t)\,ds
 =\begin{cases}
   \displaystyle\int_0^r\partial_z^m\partial_t^n
                         f_{\rm phys}(s,z,t)\,ds&k=0,\\
   \partial_r^{k-1}\partial_z^m\partial_t^n f_{\rm phys}(r,z,t)&k\ge1.
  \end{cases}
\]
The absolute values are bounded by the first term in
\eqref{mcmean:explicit_pressure_cost}. In the cutoff product all
radial derivatives act on \(\chi_m\), whereas the axial and temporal
derivatives distribute with the two displayed binomial factors.
Every derivative of its radial integral is bounded by
\(L_R\mathsf E_{0,m-i,n-j}\). This proves its second term, with
every cutoff derivative and the actual support width retained.
Angular invariance proves the stated zero entries. This construction
first computes the pressure-free \(\mathsf G\), then
\(\mathsf D,\mathsf H,\mathsf E,\mathsf\Pi\), and finally the
full force array below, so the pressure budget is a finite calculation
without circularly assuming a pressure bound.

For nonnegative arrays define their exact Leibniz convolution and
derivative shift by
\begin{equation}\label{mcmean:convolution}
 (X\star Y)_\alpha=\sum_{\beta\le\alpha}
               \binom\alpha\beta X_\beta Y_{\alpha-\beta},
 \qquad (S_\delta X)_\alpha=X_{\alpha+\delta}.
\end{equation}
Products with three factors use either parenthesization; multiplying
the binomial coefficients gives
\(\alpha!/(\beta!\gamma!(\alpha-\beta-\gamma)!)\), proving the
two parenthesizations agree. For the exact inverse-radius arrays put
\begin{equation}\label{mcmean:radius_array}
 (\mathsf R_j)_\alpha=
 \begin{cases}
 (j)_k r^{-j-k}&\ell=m=n=0,\\0&\text{otherwise},
 \end{cases}
 \quad (j)_0=1,\quad(j)_k=j(j+1)\cdots(j+k-1),\quad j=1,2.
\end{equation}
Induction gives \(\partial_r^k r^{-j}=(-1)^k(j)_kr^{-j-k}\),
which proves that \(\mathsf R_j\) is its exact absolute derivative
array. No radial derivative of an inverse-radius coefficient is
omitted in the following costs.

For \(i=r,\theta,z\) define the common transport cost
\begin{equation}\label{mcmean:transport_cost}
\begin{split}
 \mathsf T_i={}&S_{e_t}\mathsf A_i
 +\mathsf U_r\star S_{e_r}\mathsf A_i
 +\mathsf U_z\star S_{e_z}\mathsf A_i
 +\mathsf A_r\star S_{e_r}\mathsf U_i
 +\mathsf A_z\star S_{e_z}\mathsf U_i\\
 &+\mathsf R_1\star\mathsf A_\theta\star S_{e_\theta}\mathsf U_i
 +\mathsf A_r\star S_{e_r}\mathsf A_i
 +\mathsf A_z\star S_{e_z}\mathsf A_i,
\end{split}
\end{equation}
and the exact finite upper-cost arrays
\begin{equation}\label{mcmean:force_costs}
\begin{split}
 \mathsf F_r={}&\mathsf T_r
    +2\mathsf R_1\star\mathsf U_\theta\star\mathsf A_\theta
    +\mathsf R_1\star\mathsf A_\theta\star\mathsf A_\theta\\
 &+\nu_{\rm NS}\bigl(S_{2e_r}\mathsf A_r
       +\mathsf R_1\star S_{e_r}\mathsf A_r
       +S_{2e_z}\mathsf A_r+\mathsf R_2\star\mathsf A_r\bigr)
       +S_{e_r}\mathsf\Pi,\\
 \mathsf F_\theta={}&\mathsf T_\theta
    +\mathsf R_1\star\bigl(\mathsf U_r\star\mathsf A_\theta
                    +\mathsf A_r\star\mathsf U_\theta
                    +\mathsf A_r\star\mathsf A_\theta\bigr)\\
 &+\nu_{\rm NS}\bigl(S_{2e_r}\mathsf A_\theta
       +\mathsf R_1\star S_{e_r}\mathsf A_\theta
       +S_{2e_z}\mathsf A_\theta+\mathsf R_2\star\mathsf A_\theta\bigr)
       +\mathsf R_1\star S_{e_\theta}\mathsf\Pi,\\
 \mathsf F_z={}&\mathsf T_z
    +\nu_{\rm NS}\bigl(S_{2e_r}\mathsf A_z
       +\mathsf R_1\star S_{e_r}\mathsf A_z
       +S_{2e_z}\mathsf A_z\bigr)+S_{e_z}\mathsf\Pi.
\end{split}
\end{equation}
Every addition and scalar multiplication here is entrywise. For every
finite \(\alpha\), these arrays prove the full physical component
inequalities
\begin{equation}\label{mcmean:all_force_bounds}
 |\partial^\alpha\delta\mathcal R_i|\le\mathsf F_{i,\alpha},
 \qquad i=r,\theta,z.
\end{equation}
Indeed the repeated Leibniz rule gives
\(|\partial^\alpha(fg)|\le(X\star Y)_\alpha\) whenever the
respective derivative absolute values are bounded by \(X,Y\).
The three-factor version is its iterated rule, proved by the explicit
multinomial coefficient above. Apply these rules term by term to
\eqref{mcmean:force_r}--\eqref{mcmean:force_z}, using
\eqref{mcmean:velocity_bounds} for the increment, the actual
derivatives for \(U,\delta p\), and \eqref{mcmean:radius_array} for
each inverse radius. The constant viscosity passes through every
derivative. The three resulting sums are precisely
\eqref{mcmean:transport_cost}--\eqref{mcmean:force_costs}, proving
the claim at every order.

These are component derivatives in the original orthonormal basis.
For an explicit bound on the actual Cartesian vector after angular
differentiation, the moving basis gives
\[
 \left|\partial_r^k\partial_\theta^\ell\partial_z^m\partial_t^n
                 \delta\mathcal R\right|
 \le\sum_{j=0}^\ell\binom\ell j
       \bigl(\mathsf F_{r;(k,j,m,n)}
             +\mathsf F_{\theta;(k,j,m,n)}\bigr)
       +\mathsf F_{z;(k,\ell,m,n)}.
\]
To prove it, expand each angular derivative by the binomial rule
between a scalar component and its basis vector. Every derivative of
\(e_r,e_\theta\) has norm one, whereas the derivatives of \(e_z\)
of positive order vanish. The triangle inequality gives exactly the
displayed bound. Thus the component formulas also provide a bound
for the physical vector, with its angular basis differentiation.

There is also an exact recursion for every Cartesian derivative.
Let \(O_\theta\) be the orthogonal matrix with columns
\((e_r,e_\theta,e_z)\), and let
\(J(f_r,f_\theta,f_z)=(-f_\theta,f_r,0)\). Use
\(x_{\rm C}=r\cos\theta,y_{\rm C}=r\sin\theta\) for the physical
Cartesian coordinates, distinct from the original scaled radius
\(x=r/\sqrt q\). Define the operators on component vectors
\begin{equation}\label{mcmean:cartesian_recursion}
\begin{split}
 \mathcal Xf&=\cos\theta\,\partial_rf
        -\frac{\sin\theta}{r}(\partial_\theta f+Jf),\\
 \mathcal Yf&=\sin\theta\,\partial_rf
        +\frac{\cos\theta}{r}(\partial_\theta f+Jf),
 \qquad \mathcal Zf=\partial_zf.
\end{split}
\end{equation}
The scalar coordinate rules
\(\partial_{x_{\rm C}}=\cos\theta\partial_r
-r^{-1}\sin\theta\partial_\theta\) and
\(\partial_{y_{\rm C}}=\sin\theta\partial_r
+r^{-1}\cos\theta\partial_\theta\), together with
\(\partial_\theta O_\theta=O_\theta J\), give exactly
\(\partial_{x_{\rm C}}(O_\theta f)=O_\theta\mathcal Xf\),
\(\partial_{y_{\rm C}}(O_\theta f)=O_\theta\mathcal Yf\), and
\(\partial_z(O_\theta f)=O_\theta\mathcal Zf\).
Applying these identities repeatedly proves
\[
 \partial_{x_{\rm C}}^i\partial_{y_{\rm C}}^j
 \partial_z^k\partial_t^n(O_\theta f)
   =O_\theta\mathcal X^i\mathcal Y^j\mathcal Z^k\partial_t^nf.
\]
The order written is a fixed composition order and is exact because
it was obtained from commuting smooth Cartesian derivatives.
For a fully explicit majorant recursion, if \(\mathsf C_r,
\mathsf C_\theta,\mathsf C_z\) majorize the component derivative
arrays of \(f\), set \((\mathsf J C)_r=\mathsf C_\theta\),
\((\mathsf J C)_\theta=\mathsf C_r\),
\((\mathsf J C)_z=0\). Let \(\mathsf c_\alpha
=|\partial^\alpha\cos\theta|\),
\(\mathsf s_\alpha=|\partial^\alpha\sin\theta|\), whose entries
vanish unless \(k=m=n=0\) and whose remaining entries are the
actual derivatives of these elementary functions. Then one
\(\mathcal X\) application is bounded componentwise by
\[
 \mathsf c\star S_{e_r}\mathsf C_i
 +\mathsf s\star\mathsf R_1\star
             (S_{e_\theta}\mathsf C_i+(\mathsf J C)_i),
\]
and one \(\mathcal Y\) application by the same expression with
\(\mathsf c,\mathsf s\) interchanged. A \(\mathcal Z\) or time
application is the array shift \(S_{e_z}\) or \(S_{e_t}\).
These rules follow from the exact product rule
\eqref{mcmean:convolution}. Starting with \(\mathsf F\), iterating
them the indicated finite number of times, and taking the sum of
the three resulting zero-order component bounds gives a finite bound
for each stated Cartesian vector derivative. Thus all Cartesian
inverse-radius and angular-basis costs are also explicitly supplied.

All costs above are finite for the actual finite state on each compact
interior physical domain. Their dependence on the original scale,
target derivatives, inverse-radius factors, viscosity, full current
velocity, and reconstructed pressure is explicit. The exact nonlinear
defects \eqref{mcmean:nonlinear_defects} remain after this operation.
No uniform estimate over an increasing family of source stages,
no singular-time limit, and no infinite correction gain is asserted
by these finite derivative identities and costs.
